\documentclass[a4paper,5pt]{amsbook}
%%%%%%%%%%%%%%%%%%%%%%%%%%%%%%%%%%%%%%%%%%%%%%%%%%%%%%%%%%%%%%%%%%%%%

\usepackage{booktabs}
\usepackage{graphicx}
\usepackage{multicol}
\usepackage{textcomp}
\usepackage{systeme}
% \usepackage[]{float}
\usepackage{amssymb}
% \usepackage{amsfonts}
% \usepackage[]{amsmath}
% \usepackage[]{epsfig}
% \usepackage[brazil]{babel}
% \usepackage[utf8]{inputenc}
% \usepackage{verbatim}
%\usepackage[]{pstricks}
%\usepackage[notcite,notref]{showkeys}
\usepackage{subcaption}
\usepackage[inline]{enumitem}

%%%%%%%%%%%%%%%%%%%%%%%%%%%%%%%%%%%%%%%%%%%%%%%%%%%%%%%%%%%%%%

\newcommand{\sen}{\,\mbox{sen}\,}
\newcommand{\tg}{\,\mbox{tg}\,}
\newcommand{\cosec}{\,\mbox{cosec}\,}
\newcommand{\cotg}{\,\mbox{cotg}\,}
\newcommand{\ds}{\displaystyle}

%%%%%%%%%%%%%%%%%%%%%%%%%%%%%%%%%%%%%%%%%%%%%%%%%%%%%%%%%%%%%%%%%%%%%%%%

\setlength{\textwidth}{16cm} % \setlength{\topmargin}{-2cm}
% \setlength{\textheight}{25cm}
\setlength{\leftmargin}{1.2cm} \setlength{\rightmargin}{1.2cm}
\setlength{\oddsidemargin}{0cm}\setlength{\evensidemargin}{0cm}

%%%%%%%%%%%%%%%%%%%%%%%%%%%%%%%%%%%%%%%%%%%%%%%%%%%%%%%%%%%%%%%%%%%%%%%%

% \renewcommand{\baselinestretch}{1.6}
% \renewcommand{\thefootnote}{\fnsymbol{footnote}}
% \renewcommand{\theequation}{\thesection.\arabic{equation}}
% \setlength{\voffset}{-50pt}
% \numberwithin{equation}{chapter}

%%%%%%%%%%%%%%%%%%%%%%%%%%%%%%%%%%%%%%%%%%%%%%%%%%%%%%%%%%%%%%%%%%%%%%%

\begin{document}
\thispagestyle{empty}
\begin{minipage}[h]{0.14\textwidth}
	\includegraphics[scale=0.24]{../../ufgd.png}
\end{minipage}
\begin{minipage}[h]{\textwidth}
\begin{tabular}{c}
{{\bf UNIVERSIDADE FEDERAL DA GRANDE DOURADOS}}\\
{{\bf \'{A}lgebra Linear e Geometria Anal\'{\i}tica --- Lista 1}}\\
{{\bf Prof.\ Adriano Barbosa}}\\
\end{tabular}
\vspace{-0.45cm}
%
\end{minipage}

%------------------------
\vspace{1cm}

%%%%%%%%%%%%%%%%%%%%%%%%%%%%%%%%   formulario  inicio  %%%%%%%%%%%%%%%%%%%%%%%%%%%%%%%%
\begin{enumerate}
	\vspace{0.5cm}
	\item Quais das seguintes equa\c{c}\~oes s\~ao lineares em $x_1$, $x_2$ e $x_3$?
		\begin{enumerate}
			\item $x_1+5x_2-\sqrt{2}x_3=1$
			\item $x_1+3x_2+x_1x_3=2$
			\item $x_1 = -7x_2+3x_3$
			\item $x_3+x_2+x_1 = 1$
            \item $x_1^{3/5}-2x_2+x_3=4$
		\end{enumerate}

	\vspace{0.5cm}
	\item Encontre o conjunto solu\c{c}\~ao das equa\c{c}\~oes lineares.
		\begin{enumerate}
			\item $7x-5y=3$
			\item $3x-5y+4z=7$
            \item $3v-8w+2x-y+4z=0$
		\end{enumerate}

	\vspace{0.5cm}
	\item 
        \begin{enumerate}
            \item Encontre uma equa\c{c}\~ao linear nas vari\'aveis $x$ e $y$ que tem
                $x=5+2t$, $y=t$ como solu\c{c}\~ao geral.
            \item Mostre que $x=t$, $y=\frac{1}{2}t-\frac{5}{2}$ tamb\'em \'e
                solu\c{c}\~ao geral da equa\c{c}\~ao do item anterior.
        \end{enumerate}

	\vspace{0.5cm}
	\item Mostre que se o sistema
		\[\systeme{
				x+y+2z=a,
				x+z=b,
				2x+y+3z=c}\]
	tem solu\c{c}\~ao, ent\~ao as constantes $a$, $b$ e $c$ devem satisfazer $c = a+b$.

	\vspace{0.5cm}
    \item Encontre a matriz aumentada de cada um dos sistemas lineares abaixo.
        Resolva os sistemas.

		\vspace{0.3cm}
		\begin{enumerate*}
			\item
				$\systeme[x_1x_2x_3]{
					x_1 + 2x_2 - 4x_4 + x_5 = 1,
					3x_2 + x_3 - x_5 = 2,
					x_3 + 7x_4 = 1
				}$
			\hspace{1cm}
			\hspace{1cm}
			\item
				$\systeme[x_1x_2x_3]{
					2x_1 + 2x_3 = 1,
					3x_1 - x_2 + 4x_3 = 7,
					6x_1 + x_2 - x_3 = 0
				}$
		\end{enumerate*}

	\vspace{0.5cm}
	\item Encontre o sistema de equa\c{c}\~oes lineares correspondendo \`as matrizes
		aumentadas abaixo. Resolva os sistemas.

		\vspace{0.3cm}
		\begin{enumerate*}
			\item
				$\begin{bmatrix}
					2 & 0 & 0 \\
					3 & -4 & 0 \\
					0 & 1 & 1
				\end{bmatrix}$
			\hspace{1cm}
			\hspace{1cm}
			\item
				$\begin{bmatrix}
					7 & 2 & 1 & -3 & 5 \\
					1 & 2 & 4 & 0 & 1
				\end{bmatrix}$
		\end{enumerate*}

	\vspace{0.5cm}
    \item Determine se as matrizes abaixo est\~ao na forma escalonada.

		\vspace{0.3cm}
		\begin{enumerate*}
			\item
				$\begin{bmatrix}
					1 & 0 & 0 & -7 & 8 \\
					0 & 1 & 0 & 3 & 2 \\
					0 & 0 & 1 & 1 & -5
				\end{bmatrix}$
			\hspace{1cm}
			\hspace{1cm}
			\item 
				$\begin{bmatrix}
					1 & -3 & 0 & 0 \\
					0 & 0 & 1 & 0 \\
					0 & 0 & 0 & 1
				\end{bmatrix}$
		\end{enumerate*}

	\vspace{0.5cm}
	\item Resolva o sistema abaixo:
		\[\systeme[abc]{
				-2b + 3c = 1,
				3a + 6b - 3c = -2,
				6a + 6b + 3c = 5
			}\]

    \vspace{0.5cm}
    \item Dada $f(x)=ax^2+bx+c$, determine os valores de $a,b,c$ tais que $f(0)=3,
        f(-1)=2$ e $f(-3)=6$.

	\vspace{0.5cm}
	\item Para quais valores de $a$ o sistema abaixo n\~ao admite solu\c{c}\~ao?
		\[\systeme[xyz]{
				x + 2y - 3z = 4,
				3x - y + 5z = 2,
				4x + y + {(a^2-14)}z = a+2
			}\]

 	\vspace{0.5cm}
 	\item Um sistema linear \'e dito \textbf{homog\^eneo} se os coeficientes
 		independentes s\~ao todos zero, ou seja, a \'ultima coluna da matriz
 		aumentada do sistema \'e nula. \'{E} poss\'{\i}vel que um sistema linear
 		homog\^eneo n\~ao possua solu\c{c}\~ao?
\end{enumerate}

\end{document}
